\documentclass[a4paper]{article}

\usepackage[utf8]{inputenc}
\usepackage[T1]{fontenc}
\usepackage[italian]{babel}
\usepackage{titling}
\usepackage{geometry}
\usepackage{amsmath}

\makeatletter
\newcommand{\subtitle}[1]{\gdef\@subtitle{#1}}
\newcommand{\@subtitle}{}
\pretitle{
  \begin{center}
  \LARGE
}
\posttitle{
  \par\vskip0.5em
  \large\@subtitle{}
  \end{center}
}
\makeatother

\title{Tutorato 08}
\subtitle{Boss fight}
\author{Prigione Luca}
\date{04/12/XXXX}

\begin{document}
\pagestyle{empty}

\maketitle

\section*{Compleanni}
Dopo giorni di prove, allenamenti e rivelazioni, Lele raggiunse finalmente il sentiero che conduceva al castello dove, secondo ogni leggenda, viveva l’entità più temuta del regno digitale: il Boss.
La fortezza era enorme e silenziosa. Gab, Ballocco e Tuffu lo avevano aiutato fino a quel momento, ma decisero di non avventurarsi oltre il ponte levatoio.\\
Tuffu annuì, seria:

«Ricorda: prima di affrontarlo dovrai dimostrare di saper controllare le costruzioni più complesse. Il castello riconosce solo chi sa dare forma ai dati.»\\
Per riuscire a farvi accettare dal castello e poichè Lele è un cavaliere gentile, decidete di costruire un calendario utilizzando delle liste in modo da segnarvi le date di compleanno dei personaggi che avete incontrato lungo il percorso.\\
Dovrete quindi implementare una funzione \texttt{crea\_calendario}, una funzione \texttt{add\_compleanno}, una funzione \texttt{stampa\_calendario} e una funzione \texttt{elimina\_calendario} affinché il calendario funzioni correttamente.

\section*{La sala del trono}
Il salone era immenso. Lampadari di cristallo fluttuavano sopra il vuoto, e in fondo, su un trono spezzato, sedeva un gigante d’ombra e luce: il Boss.\\
Non aveva titolo, non aveva identità: era semplicemente il Boss.\\
Guardiano dei dati, distruttore delle strutture mal progettate.

«Sei arrivato fin qui…» rimbombò la sua voce.

«Mostrami se meriti di ristrutturare questo mondo.»\\
Il colosso si alzò, e il suo corpo si divise in tre forme, ognuna rivelatrice di un aspetto del caos delle liste.

\section*{La forma disordinata}
La prima forma era puro caos: elementi sparsi, ripetuti, privi di ordine e logica.\\
Per riuscire a sconfiggerla, dovete creare un programma che gestisca una lista non ordinata.
Per farla funzionare, dovrete implementare le seguenti funzioni:
\begin{itemize}
    \item \texttt{inserisci (lista, valore)}: inserisce un nuovo nodo con il valore specificato in testa alla lista.
    \item \texttt{rimuovi\_duplicati (lista)}: rimuove tutti i valori all'interno della lista che sono duplicati.
    \item \texttt{stampa\_lista (lista)}: stampa tutti i valori presenti nella lista.
\end{itemize}

\section*{La forma specchiata}
La seconda forma si componeva di specchi. Ogni suo movimento generava una copia riflessa che si muoveva al contrario, come se la realtà stessa oscillasse tra due direzioni opposte.\\
Per superarla, dovete creare un programma che gestisca una lista doppiamente concatenata di numeri. Per farla funzionare, dovrete implementare le seguenti funzioni:
\begin{itemize}
    \item \texttt{inserisci\_in\_posizione (lista, posizione, valore)}: inserisce un nuovo nodo con il valore specificato nella posizione corretta.
    \item \texttt{rimuovi\_in\_posizione (lista, posizione)}: rimuove il nodo alla posizione corretta.
    \item \texttt{rimuovi\_valore (lista, valore)}: rimuove il primo nodo che contiene il valore specificato dalla lista.
    \item \texttt{stampa\_lista (lista)}: stampa tutti i valori presenti nella lista in entrambi i versi.
\end{itemize}

\section*{La forma finale}
La terza forma era un anello di energia: un ciclo infinito, senza inizio né fine.

«VEDIAMO SE PUOI INTERROMPERE L’ETERNO…»\\
Per sconfiggere definitivamente il Boss, dovete scrivere un programma che costruisca una lista circolare di interi, implementando le seguenti funzioni ricorsive:
\begin{itemize}
    \item \texttt{inserisci (lista, posizione)}: inserisce al fondo un nuovo nodo con il valore specificato.
    \item \texttt{lunghezza (lista)}: percorre la lista e calcola il numero di nodi presenti.
    \item \texttt{remove\_max (lista)}: rimuove il nodo contenente il valore massimo, mantenendo la circolarità.
    \item \texttt{stampa (lista)}: stampa tutti i valori presenti nella lista, partendo dal nodo iniziale e tornando al nodo iniziale.
\end{itemize}
P.S.: Se non avete utilizzato le liste circolari in precedenza, potete pensare di utilizzare gli array circolari.

\noindent\rule{\linewidth}{0.4pt}
\section*{Esercizi aggiuntivi}
Se avete già finito e volete qualche esercizio in più chiedete al Tutor, così potete provare a testare il vostro livello.

\end{document}